\section{Latar Belakang}

Mejiro Express ingin menguji coba layanan logistik darat dengan satu truk peti kemas.
Setiap pengiriman perlu mengatur paket mana yang dimuat, di mana paket diletakkan,
dan bagaimana paket diputar agar total nilai (\textit{value}) paket yang terangkut
sebesar mungkin, tanpa melanggar batasan fisik peti kemas. Persoalan ini merupakan
varian persoalan \textit{3D knapsack} / \textit{3D bin packing} yang tergolong NP-hard,
sehingga pencarian solusi eksak tidak praktis untuk ukuran instans yang realistis.
Tugas ini menyelesaikannya dengan algoritma \textit{local search}: Hill-Climbing,
Simulated Annealing, dan Genetic Algorithm~\cite{asisten2026spek}.

\section{Variabel Persoalan}

Setiap paket berbentuk balok dengan atribut berikut.

\begin{table}[H]
	\centering
	\caption{Atribut paket dan truk}
	\label{tab:atribut}
	\begin{tabular}{l l p{8cm}}
		\toprule
		\textbf{Entitas} & \textbf{Atribut} & \textbf{Keterangan} \\
		\midrule
		Paket & \texttt{id} & Identitas paket \\
		& \texttt{w, l, h} & Lebar, panjang, dan tinggi paket \\
		& \texttt{orientation} & Sumbu truk yang sejajar dengan sisi \texttt{w}, \texttt{l}, \texttt{h}; ada enam kemungkinan \\
		& \texttt{value} & Nilai paket \\
		& \texttt{weight} & Berat paket \\
		& \texttt{fragile} & Paket pecah-belah atau tidak \\
		& \texttt{eta} & Perkiraan sisa hari hingga paket harus tiba \\
		\midrule
		Truk & \texttt{w, l, h} & Dimensi peti kemas \\
		& \texttt{max\_capacity} & Berat maksimum yang dapat diangkut \\
		\bottomrule
	\end{tabular}
\end{table}

\section{Fungsi Objektif dan Batasan}

Fungsi objektif yang dimaksimalkan adalah total \textit{value} paket yang berada di
dalam peti kemas:
\begin{equation}
	f(s) = \sum_{i \,:\, \text{paket } i \text{ di dalam truk}} \text{value}_i .
\end{equation}
Paket yang berada di luar truk tidak menyumbang nilai. Sebuah konfigurasi dianggap
valid apabila memenuhi seluruh batasan berikut.

\begin{enumerate}
	\item \textbf{Tidak menembus dinding.} Seluruh bagian paket di dalam truk berada di dalam peti kemas.
	\item \textbf{Tidak saling bertumpuk ruang.} Dua paket tidak boleh menempati ruang yang sama.
	\item \textbf{Tidak melayang.} Paket dengan sudut kiri-bawah-belakang $(x_i,y_i,z_i)$ dan alas
	$B_i=\{(x,y,z_i)\mid x_i\le x<x_i+w_i,\ y_i\le y<y_i+d_i\}$ harus memiliki minimal satu titik
	alas yang ditopang, yaitu $z_i=0$ (lantai truk) atau tepat berada di atas permukaan atas paket lain.
	\item \textbf{Barang pecah-belah tidak menjadi penopang.} Tidak boleh ada paket yang bertumpu di atas paket \textit{fragile}.
	\item \textbf{Kapasitas.} Total berat paket di dalam truk tidak melebihi \texttt{max\_capacity}.
\end{enumerate}

Koordinat berupa bilangan bulat. State awal dibangkitkan secara acak. Aksi tetangga yang
diperbolehkan adalah (1) menukar dua paket, baik di dalam maupun di luar truk,
(2) memindahkan paket ke koordinat lain, dan (3) merotasi paket $90^\circ$ terhadap
sumbu $x$, $y$, atau $z$.

\section{Instans Uji}

Tiga instans dibangkitkan dengan \texttt{src/generate\_instance.py} (\textit{seed} 42) dan
disimpan di \texttt{data/instances/}. Kolom ``batas atas'' adalah nilai optimum
\textit{knapsack} 0/1 dengan batasan berat dan volume yang mengabaikan geometri
(Subbab~\ref{sec:bound}); nilai ini tidak lebih kecil daripada optimum global persoalan.

\begin{table}[H]
	\centering
	\caption{Instans uji}
	\label{tab:instans}
	\begin{tabular}{l c c c c}
		\toprule
		\textbf{Instans} & \textbf{Truk} ($w\times l\times h$) & \textbf{Kapasitas} & \textbf{Jumlah paket} & \textbf{Batas atas} \\
		\midrule
		small  & $6\times8\times5$    & 150  & 15  & 734 \\
		medium & $10\times20\times5$  & 500  & 55  & 2384 \\
		large  & $12\times24\times8$  & 1200 & 120 & 5633 \\
		\bottomrule
	\end{tabular}
\end{table}
